\subsection{Soundness of Additional Rewrite Schemas}
\label{app:axis-extensions}
This appendix proves soundness of specific rewrite schemas beyond the
complete constructor's fragment. Each schema has explicit conditions under
which its generated attributes preserve the written rewrite. Supplied RHS
fields are preserved only when they are certified equal to the constructed
fields. The result covers the listed schemas, not arbitrary compositions
or completeness of inference and guard checking.

\paperparagraph{Coordinate interpretation.}
We compare the input entries read at each output coordinate. Write
$B_{d,e}X$ for insertion-only broadcast: the distinct positions $d$ locate the original
input dimensions, and $e$ lists inserted extents in output order. Thus
$(B_{d,e}X)[z]=X[z[d]]$. A transpose $T_\phi X$ has shape
$\gath(\operatorname{shape}(X),\phi)$ and reads $X[z[\phi^{-1}]]$.
A reversal $R_AX$ replaces coordinate $z_i$ by $n_i-1-z_i$ for $i\in A$,
where $n_i$ is that dimension's extent. A reduction $\Sigma_AX$ sums over
positions $A$ and preserves the order of the other dimensions. All source
selectors below are distinct and in range.

\begin{lemma}[Sound rewrite schemas]
\label{lem:certified-extensions}
Under the canonical declaration and dtype contracts, each constructor
below produces a legal replacement with identical ordered shape and values
whenever its stated source condition holds. This holds for arbitrary finite
ranks and nonnegative extents under exact algebraic semantics, sufficient
resources, and representable metadata.
\end{lemma}
\begin{proof}
We group the schemas by the argument needed to establish their correctness.
For each group, we give the construction, its conditions, and the reason
it preserves input reads or finite sums. We then collect the legality,
shape, and dtype arguments.

\paperparagraph{Step 1: Simplify views without changing input reads.}
The first task is to remove redundant view operations while preserving
which input entry each output coordinate reads.
A chain of reversals becomes one reversal on the symmetric difference
of its axis sets. Reversing a coordinate twice restores it, and reversals
on different coordinates commute. The remaining reversal can be removed
only when its axis set is empty.

For nested broadcasts $B_{b,f}(B_{d,e}X)$, require
$b=[0,\ldots,r+|e|-1]$, where $r$ is the input rank. The outer broadcast
then appends its inserted dimensions. The replacement embedding is
$\gath(b,d)$, and its inserted extents are $e\catseq f$. This preserves
both the original input coordinates and the order of inserted dimensions.

For the singleton-axis view, require
$\operatorname{shape}(X)=[1]\catseq t$ and $d=[0,\ldots,r-1]$,
where $r=\operatorname{rank}(X)$.
Remove input dimension zero and broadcast the remaining tensor with
embedding $[1,\ldots,r-1]$ and inserted extents $[1]\catseq e$.
Both results have shape $[1]\catseq t\catseq e$ and read the same entry:
the removed input coordinate can only be zero. This argument also applies
to nonnumeric elements; it does not provide a general reshape rule.

\paperparagraph{Step 2: Translate operations through broadcast.}
Broadcast introduces output coordinates that do not affect the input read.
We must distinguish these inserted dimensions from the embedded original
dimensions when moving another operation across broadcast.

For $\star\in\{+,\times\}$, matching embeddings and inserted extents give
$(B_{d,e}X)\star(B_{d,e}Y)=B_{d,e}(X\star Y)$.
Source pointwise legality and the common injective embedding imply equal
input shapes. At coordinate $z$, both sides compute
$X[z[d]]\star Y[z[d]]$.

For reversal, set $Q=\{j\mid d_j\in A\}$. Then
$R_A(B_{d,e}X)=B_{d,e}(R_QX)$: reversing an inserted dimension cannot
change the input read, while reversing an embedded dimension reverses its
corresponding input coordinate. The input reversal can be omitted only
when $Q=\varnothing$; this includes scalar broadcast with $d=[]$.

For transpose, let $t=\operatorname{shape}(B_{d,e}X)$ and $n=|t|$.
Construct
\[
 d'=\gath(\phi^{-1},d),\qquad
 e'=\gath(\gath(t,\phi),\ocomp([0,\ldots,n-1],d')).
\]
Original input dimension $j$ moves from position $d_j$ to
$\phi^{-1}[d_j]$. The list $d'$ records those positions, and $e'$ records
the other extents in the transposed output order. Hence $B_{d',e'}X$
has the same shape and reads the same input coordinates as
$T_\phi(B_{d,e}X)$. For scalar broadcast, this reduces to $d'=[]$ and
$e'=\gath(e,\phi)$. A broadcast inserting no dimensions is itself the
transpose $T_{d^{-1}}X$.

For reduction, require $A\subseteq\operatorname{im}(d)$: only embedded
original dimensions may be reduced. Set $Q=\{j\mid d_j\in A\}$ and, for
each $j\notin Q$, retain the embedding entry
$d_j-|\{a\in A\mid a<d_j\}|$. These entries form $d'$ in input order,
giving $\Sigma_A(B_{d,e}X)=B_{d',e}(\Sigma_QX)$.
Both sides sum over the same input coordinates. The subtraction shifts
surviving positions after the reduced dimensions are removed. Reducing an
inserted dimension is excluded because it would add a multiplicity factor.

\paperparagraph{Step 3: Preserve scalar algebra at each coordinate.}
These schemas move scalar operations without changing which values are
combined. Reversal distributes over addition because both sides add the
same two reversed input entries. A broadcast scalar is constant under
reversal and can be copied unchanged. A chain of numeric negations depends
only on its parity, and negation commutes with finite summation, including
an empty sum.

For scalar $c$, bilinearity gives
$\operatorname{dot}(cX,Y)=c\operatorname{dot}(X,Y)$, and similarly on the
right. Copy the dot's axis fields and broadcast $c$ to the original dot's
output shape, computed from its ordered batch and free positions.
For nested products, the identity
\[
 (X\cdot B_{[],\operatorname{shape}(X)}c)\cdot B_{d,e}Y
 =X\cdot B_{d,e}(Y\cdot B_{[],\operatorname{shape}(Y)}c)
\]
uses source pointwise legality to equate the outer shapes. At each
coordinate, it is associativity and commutativity of the same three scalar
factors. Identical operation trees can also copy their source attributes:
induction over their matching nodes preserves each operation and result.

\paperparagraph{Step 4: Preserve contraction partners.}
When moving an operation between dot operands, the important requirement
is to keep the same paired coordinates.
To move reversal from one operand to the other, require every reversed
position in $A$ to belong to that operand's contracting list $k$.
If $h$ is the other contracting list, reverse positions
$\gath(h,\{j\mid k_j\in A\})$ there and copy all dot axis fields.
Source legality gives equal extents for each pair. Replacing a selected
summation index $t$ by $n-1-t$ is a bijection, so the finite sum is unchanged;
batch and free coordinates are unaffected.

Two dot conversions follow directly from the same pairing interpretation.
With no batch axes and both contracting lists equal to their full input
position ranges, the operands have identical ordered shapes. The dot is
therefore their pointwise product followed by reduction of all dimensions.
With all four axis lists empty, the dot is instead an outer product.
Broadcast the left operand into the output prefix and the right operand
into the output suffix, inserting the other operand's shape in each case,
then multiply. This preserves the dot's left-free/right-free output order.

\paperparagraph{Step 5: Move broadcast-created contractions into a reduction.}
Here an inserted broadcast dimension participates in a dot contraction.
We must preserve its summation even though the original unbroadcast input
does not contain that dimension.
Let $d$ embed $X$ into the left operand, and let $I$ be the inserted
positions. Require every position in $I$ to belong to the left contracting
list. Reduce $Y$ over the right contracting positions paired with $I$,
then remove those pairs from the dot. Map the remaining left batch and
contracting positions back through $d$. Shift each surviving right position
to account for the preceding dimensions removed by the reduction.

Require left batch positions to be embedded original positions and the
mapped left free positions to preserve the source's free output order.
The removed right positions are distinct and disjoint from retained batch
and contracting positions. These conditions make the new reduction and
dot legal and preserve batch/left-free/right-free output order.
At each fixed output coordinate, the value identity is
\[
 \sum_{a,b}X[a]Y[a,b]=\sum_a X[a]\left(\sum_bY[a,b]\right),
\]
where $b$ ranges over only the broadcast-created contracting dimensions.
Finite-sum distribution proves equality, including empty summation domains.

\paperparagraph{Step 6: Establish the full replacement guarantee.}
The preceding arguments identify the same input reads or equal finite
sums for each schema. Their stated conditions also ensure distinct,
in-range selectors, disjoint batch and contracting roles, and the required
output order. Broadcast constructors preserve the ordered output shape;
reductions remove exactly the specified dimensions; dot constructors keep
the declared batch and free-axis order.

Views preserve dtype, and the numeric operations preserve the compatible
dtypes required by the source contracts. Thus source legality establishes
the replacement's dtype requirements. A supplied field proved equal to
the constructed field under the effective applicability condition gives
the same attributes.
For a zero extent, there are no coordinates to read on that dimension,
and reduction over it is an empty sum. For extent one, reversal is the
identity and the singleton-axis construction reads coordinate zero.
The arguments therefore apply to nonnegative extents without the
completeness theorem's minimum-extent assumption. This proves soundness
of the listed schemas while leaving that theorem's fragment unchanged.
\end{proof}
